\documentclass[8pt,letterpaper]{article}

\usepackage[utf8]{inputenc}
\usepackage{amsmath}
\usepackage{amssymb}
\usepackage[margin=0.5cm]{geometry}
\usepackage{multicol}
\usepackage[compact]{titlesec}
\usepackage{enumitem}

% Ajustes para títulos pequeños y compactos
\titleformat{\section}[runin]{\normalfont\bfseries}{\thesection.}{0.5em}{}[.]
\titlespacing*{\section}{0pt}{0.5ex plus 0.5ex minus 0.2ex}{0.5ex plus 0.2ex}
\titleformat{\subsection}[runin]{\normalfont\itshape}{\thesubsection.}{0.5em}{}[.]
\titlespacing*{\subsection}{0pt}{0.5ex plus 0.5ex minus 0.2ex}{0.5ex plus 0.2ex}

\pagestyle{empty} % Sin números de página

\begin{document}

\begin{multicols*}{3}

% =============================================================================
% PÁGINA 1: LÓGICA Y DESIGUALDADES
% =============================================================================
\section*{LÓGICA Y DESIGUALDADES}

\subsection*{Demostración por el Absurdo}
Para probar que $p$ es irracional (ej. $\sqrt{2}$):
\begin{enumerate}[label=\arabic*.]
    \item Asumir que $p$ es racional: $p = a/b$, donde $a, b \in \mathbb{Z}$, $b \neq 0$, y $\text{mcd}(a,b)=1$.
    \item Manipular la ecuación para llegar a una contradicción (ej. $a$ y $b$ son ambos pares).
    \item Concluir que la suposición inicial es falsa, por lo tanto $p$ es irracional.
\end{enumerate}

\subsection*{Inecuaciones con Valor Absoluto}
\begin{itemize}
    \item $|A| \le k \iff -k \le A \le k$ (si $k \ge 0$)
    \item $|A| \ge k \iff A \ge k \text{ o } A \le -k$ (si $k \ge 0$)
    \item \textbf{Método de Puntos Críticos:}
    \begin{enumerate}[label=\arabic*.]
        \item Hallar raíces del numerador y denominador (para racionales) o valores que anulan los valores absolutos.
        \item Ubicar puntos en la recta numérica y formar intervalos.
        \item Evaluar el signo de la expresión en cada intervalo.
        \item Seleccionar intervalos que satisfacen la inecuación.
    \end{enumerate}
\end{itemize}

\subsection*{Axioma del Supremo}
\begin{itemize}
    \item \textbf{Cota Superior:} $M$ es cota superior de $S$ si $\forall x \in S, x \le M$.
    \item \textbf{Cota Inferior:} $m$ es cota inferior de $S$ si $\forall x \in S, x \ge m$.
    \item \textbf{Supremo (sup):} La menor de las cotas superiores. Si $\sup S \in S$, es el \textbf{Máximo}.
    \item \textbf{Ínfimo (inf):} La mayor de las cotas inferiores. Si $\inf S \in S$, es el \textbf{Mínimo}.
\end{itemize}

\end{multicols*}
\newpage
\begin{multicols*}{3}

% =============================================================================
% PÁGINA 2: GEOMETRÍA ANALÍTICA Y CÓNICAS
% =============================================================================
\section*{GEOMETRÍA ANALÍTICA Y CÓNICAS}

\subsection*{La Recta}
\begin{itemize}
    \item \textbf{Pendiente $m$:} $m = \frac{y_2 - y_1}{x_2 - x_1}$.
    \item \textbf{Paralelismo:} $m_1 = m_2$.
    \item \textbf{Perpendicularidad:} $m_1 \cdot m_2 = -1$.
\end{itemize}

\subsection*{Circunferencia}
\begin{itemize}
    \item \textbf{Ecuación Canónica:} $(x-h)^2 + (y-k)^2 = r^2$.
    \item \textbf{Centro:} $(h, k)$. \textbf{Radio:} $r$.
    \item \textbf{Intersección de dos circunferencias:}
    \begin{enumerate}[label=\arabic*.]
        \item Restar las ecuaciones generales para obtener una recta.
        \item Sustituir la recta en una de las ecuaciones de la circunferencia.
        \item Resolver la ecuación cuadrática resultante para hallar los puntos de intersección.
    \end{enumerate}
\end{itemize}

\subsection*{Cónicas}
\begin{itemize}
    \item \textbf{Completación de Cuadrados:}
    $ax^2+bx = a\left(x+\frac{b}{2a}\right)^2 - \frac{b^2}{4a}$.
    \item \textbf{Parábola:}
    \begin{itemize}
        \item $(x-h)^2 = 4p(y-k)$ (abre hacia arriba/abajo)
        \item $(y-k)^2 = 4p(x-h)$ (abre hacia derecha/izquierda)
        \item \textbf{Vértice:} $(h,k)$. \textbf{Foco:} $(h, k+p)$ o $(h+p, k)$. \textbf{Directriz:} $y=k-p$ o $x=h-p$.
    \end{itemize}
    \item \textbf{Elipse:}
    \begin{itemize}
        \item $\frac{(x-h)^2}{a^2} + \frac{(y-k)^2}{b^2} = 1$ (eje mayor horizontal)
        \item $\frac{(x-h)^2}{b^2} + \frac{(y-k)^2}{a^2} = 1$ (eje mayor vertical)
        \item $c^2 = a^2 - b^2$. \textbf{Centro:} $(h,k)$. \textbf{Focos:} $(h \pm c, k)$ o $(h, k \pm c)$. \textbf{Vértices:} $(h \pm a, k)$ o $(h, k \pm a)$.
    \end{itemize}
    \item \textbf{Hipérbola:}
    \begin{itemize}
        \item $\frac{(x-h)^2}{a^2} - \frac{(y-k)^2}{b^2} = 1$ (eje transverso horizontal)
        \item $\frac{(y-k)^2}{a^2} - \frac{(x-h)^2}{b^2} = 1$ (eje transverso vertical)
        \item $c^2 = a^2 + b^2$. \textbf{Centro:} $(h,k)$. \textbf{Focos:} $(h \pm c, k)$ o $(h, k \pm c)$. \textbf{Vértices:} $(h \pm a, k)$ o $(h, k \pm a)$.
        \item \textbf{Asíntotas:} $y-k = \pm \frac{b}{a}(x-h)$ (horizontal) o $y-k = \pm \frac{a}{b}(x-h)$ (vertical).
    \end{itemize}
\end{itemize}

\end{multicols*}
\newpage
\begin{multicols*}{3}

% =============================================================================
% PÁGINA 3: TRIGONOMETRÍA Y LÍMITES
% =============================================================================
\section*{TRIGONOMETRÍA Y LÍMITES}

\subsection*{Trigonometría}
\begin{itemize}
    \item \textbf{Identidad de la suma de senos:}
    $\sin^2\alpha + \sin^2\beta + \sin^2\gamma = 4\sin\alpha\sin\beta\sin\gamma$ (para $\alpha+\beta+\gamma=\pi$).
    \item \textbf{Ley de Senos:}
    $\frac{a}{\sin A} = \frac{b}{\sin B} = \frac{c}{\sin C}$. Útil para problemas de funiculares/cerros.
\end{itemize}

\subsection*{Límites Notables}
\begin{itemize}
    \item $\lim_{x\to 0} \frac{\sin x}{x} = 1$
    \item $\lim_{x\to 0} \frac{1 - \cos(kx)}{x^2} = \frac{k^2}{2}$
\end{itemize}

\subsection*{Técnicas de Límites}
\begin{itemize}
    \item \textbf{Racionalización:} Multiplicar por el conjugado para eliminar raíces.
    \item \textbf{División por la mayor potencia (para $x \to \infty$):}
    Dividir numerador y denominador por la mayor potencia de $x$.
    \item \textbf{Teorema del Sándwich (o de Compresión):}
    Si $g(x) \le f(x) \le h(x)$ para todo $x$ en un intervalo abierto que contiene a $c$ (excepto posiblemente en $c$), y $\lim_{x\to c} g(x) = L = \lim_{x\to c} h(x)$, entonces $\lim_{x\to c} f(x) = L$.
\end{itemize}

\subsection*{Continuidad en un Punto $c$}
Una función $f(x)$ es continua en $c$ si se cumplen las 3 condiciones:
\begin{enumerate}[label=\arabic*.]
    \item $f(c)$ está definida.
    \item $\lim_{x\to c} f(x)$ existe.
    \item $\lim_{x\to c} f(x) = f(c)$.
\end{enumerate}
\textbf{Para funciones por partes:} Igualar los límites laterales en los puntos de quiebre para hallar parámetros $a, b$.

\end{multicols*}
\newpage
\begin{multicols*}{3}

% =============================================================================
% PÁGINA 4: DERIVADAS Y ASÍNTOTAS
% =============================================================================
\section*{DERIVADAS Y ASÍNTOTAS}

\subsection*{Análisis de Asíntotas}
\begin{itemize}
    \item \textbf{Verticales:} $x=c$ si $\lim_{x\to c^\pm} f(x) = \pm \infty$. Ocurre cuando el denominador es cero y el numerador no.
    \item \textbf{Horizontales:} $y=L$ si $\lim_{x\to \pm \infty} f(x) = L$.
    \item \textbf{Oblicuas:} $y=mx+b$.
    \begin{itemize}
        \item $m = \lim_{x\to \pm \infty} \frac{f(x)}{x}$
        \item $b = \lim_{x\to \pm \infty} (f(x) - mx)$
    \end{itemize}
\end{itemize}

\subsection*{Derivadas}
\begin{itemize}
    \item \textbf{Regla del Producto:} $(uv)' = u'v + uv'$.
    \item \textbf{Regla del Cociente:} $\left(\frac{u}{v}\right)' = \frac{u'v - uv'}{v^2}$.
    \item \textbf{Regla de la Cadena:} $(f(g(x)))' = f'(g(x)) \cdot g'(x)$.
    \item \textbf{Derivación Implícita:} Derivar ambos lados de la ecuación con respecto a $x$, tratando $y$ como función de $x$ ($y'$).
    \item \textbf{Derivadas de Orden Superior:} $y'', y'''$, etc.
\end{itemize}

\subsection*{Recta Tangente}
La ecuación de la recta tangente a $f(x)$ en el punto $(x_0, y_0)$ es:
$y - y_0 = f'(x_0)(x - x_0)$.

\end{multicols*}

\end{document}
